\documentclass[12pt,letterpaper,oneside]{extreport}

% =========================================
% PAQUETES
% =========================================
\usepackage[utf8]{inputenc}
\usepackage[T1]{fontenc}
\usepackage[spanish]{babel}

\usepackage{geometry}
\geometry{
    letterpaper,
    left=3cm,
    right=3cm,
    top=2.5cm,
    bottom=2.5cm
}

\usepackage{setspace}
\onehalfspacing

\usepackage{graphicx}
\usepackage{hyperref}
\usepackage{xcolor}
\usepackage{listings}
\usepackage{longtable}
\usepackage{booktabs}
\usepackage{float}
\usepackage{enumitem}
\usepackage{titlesec}
\usepackage{parskip}

% TikZ for ER diagrams
\usepackage{tikz}
\usetikzlibrary{arrows.meta,positioning,shapes.multipart,calc}

% Styles for ER diagram
	\tikzset{
    entity/.style={draw, thick, rectangle, rounded corners, fill=gray!10, inner sep=6pt, minimum width=3.5cm},
    attr/.style={draw, thin, ellipse, fill=white!90, inner sep=2pt},
    rel/.style={draw, diamond, aspect=1.5, inner sep=2pt, fill=white!90},
    pk/.style={font=\bfseries},
    fk/.style={font=\itshape},
    line/.style={-Latex, >=Stealth, thick}
}

% =========================================
% ELIMINAR ESPACIOS FEOS DE CHAPTER
% =========================================
\titlespacing*{\chapter}{0pt}{-40pt}{20pt}

% =========================================
% CONFIGURACIÓN DE TÍTULOS
% =========================================
\titleformat{\chapter}
{\Huge\bfseries}
{\thechapter.}
{1em}
{}

\titleformat{\section}
{\LARGE\bfseries}
{\thesection}
{1em}
{}

\titleformat{\subsection}
{\Large\bfseries}
{\thesubsection}
{1em}
{}

% =========================================
% CONFIGURACIÓN DE CÓDIGO
% =========================================
\lstset{
    basicstyle=\ttfamily\normalsize,
    breaklines=true,
    frame=single,
    backgroundcolor=\color{gray!10}
}

% =========================================
% DOCUMENTO
% =========================================
\begin{document}

% =========================================
% PORTADA
% =========================================
\begin{titlepage}

\centering

\vspace*{1.5cm}

{\Huge \textbf{Universidad Nacional Autónoma de México}}\\[0.8cm]

{\LARGE Facultad de Estudios Superiores Aragón}\\[1.8cm]

{\Huge \textbf{DOCUMENTO INTEGRAL DEL PROYECTO}}\\[1.2cm]

{\Huge Sistema de Gestión de Incidencias}\\[2cm]

\begin{flushleft}

\Large

\textbf{Asignatura:} Desarrollo Web\\[0.5cm]

\textbf{Proyecto:} Sistema de Incidencias\\[0.5cm]

\textbf{Integrantes:} Equipo de Desarrollo\\[0.5cm]

\textbf{Institución:} FES Aragón\\[0.5cm]

\textbf{Fecha de entrega:} \today

\end{flushleft}

\vfill

\end{titlepage}

% =========================================
% ÍNDICE
% =========================================
\tableofcontents
\newpage

% =========================================
% RESUMEN EJECUTIVO
% =========================================
\chapter{Resumen Ejecutivo}

El presente proyecto denominado ``Sistema de Gestión de Incidencias'' fue desarrollado con la finalidad de optimizar y modernizar el proceso de administración de incidencias dentro de la Facultad de Estudios Superiores Aragón.

Actualmente, muchos reportes relacionados con mantenimiento, infraestructura, aulas, laboratorios y equipos tecnológicos son realizados mediante métodos manuales o medios informales, provocando pérdida de información, retrasos administrativos y falta de seguimiento adecuado.

Para solucionar esta problemática se diseñó una plataforma web que permite centralizar toda la información relacionada con incidencias institucionales.

El sistema permite:

\begin{itemize}
    \item Registro e inicio de sesión de usuarios.
    
    \item Administración de incidencias.
    
    \item Seguimiento de reportes.
    
    \item Gestión de estados.
    
    \item Visualización de estadísticas.
    
    \item Administración de usuarios.
\end{itemize}

Las tecnologías utilizadas incluyen PHP 8, MySQL 8, Docker (Docker Compose), HTML5, CSS3, JavaScript y GitHub.
Adicionalmente el proyecto incorpora librerías del frontend para visualización y exportación: Chart.js para gráficas y html2canvas + jsPDF para exportar el dashboard a PDF.

El despliegue está orquestado con \texttt{docker-compose.yml} y define dos servicios principales: el servicio web (contenedor llamado \texttt{fes\_web}) y el servicio de base de datos (contenedor llamado \texttt{fes\_mysql}).
Por defecto se exponen los puertos 8080 (host -> contenedor Apache 80) y 3307 (host -> contenedor MySQL 3306). La base de datos creada por defecto se llama \texttt{fesaragon} y la contraseña root usada en el entorno de desarrollo es \texttt{root} (ver sección de seguridad para recomendaciones de producción).

Credenciales de ejemplo incluidas en \texttt{database/seed.sql} (desarrollo): usuario administrador \texttt{admin@aragon.unam.mx} con contraseña \texttt{Admin2026}. El seed original contiene contraseñas en texto para facilitar pruebas; el sistema implementa migración a \texttt{password\_hash} y se recomienda reemplazar estas entradas por hashes antes de producción.

El sistema busca mejorar la eficiencia administrativa y contribuir a la transformación digital dentro de la institución.

% =========================================
% INTRODUCCIÓN
% =========================================
\chapter{Introducción}

La transformación digital se ha convertido en una necesidad importante para las instituciones educativas modernas.

Actualmente, muchos procesos administrativos continúan realizándose manualmente, generando problemas de organización y retrasos operativos.

Dentro de la FES Aragón existen constantemente incidencias relacionadas con infraestructura, laboratorios, centros de cómputo y servicios tecnológicos que requieren seguimiento continuo.

Sin embargo, la ausencia de una plataforma centralizada dificulta el control adecuado de la información.

Debido a esta problemática surge la necesidad de desarrollar un sistema web capaz de centralizar y administrar incidencias de manera eficiente.

El proyecto busca optimizar procesos administrativos y mejorar la comunicación entre usuarios y administradores.

% =========================================
% PLANTEAMIENTO DEL PROBLEMA
% =========================================
\chapter{Planteamiento del Problema}

Actualmente, la gestión de incidencias presenta múltiples limitaciones operativas.

Muchos reportes relacionados con mantenimiento o fallas tecnológicas se realizan mediante métodos informales.

Entre los principales problemas detectados se encuentran:

\begin{itemize}
    \item Pérdida de información.
    
    \item Duplicidad de reportes.
    
    \item Falta de seguimiento.
    
    \item Ausencia de estadísticas.
    
    \item Procesos administrativos lentos.
\end{itemize}

Estas limitaciones afectan directamente la eficiencia administrativa y la calidad de atención ofrecida a estudiantes y personal académico.

Por ello se propone desarrollar una plataforma tecnológica que permita automatizar el proceso de administración de incidencias.

% =========================================
% OBJETIVOS
% =========================================
\chapter{Objetivos}

\section{Objetivo General}

Desarrollar un sistema web de gestión de incidencias para optimizar el registro y seguimiento de reportes dentro de la FES Aragón.

\section{Objetivos Específicos}

\begin{enumerate}
    \item Implementar autenticación segura.
    
    \item Desarrollar un módulo CRUD de incidencias.
    
    \item Centralizar información administrativa.
    
    \item Implementar estadísticas administrativas.
    
    \item Facilitar el seguimiento de incidencias.
    
    \item Mejorar la organización institucional.
\end{enumerate}

% =========================================
% MARCO TEÓRICO
% =========================================
\chapter{Marco Teórico}

\section{Sistemas de Información}

Los sistemas de información permiten recopilar, procesar y almacenar datos con el objetivo de apoyar procesos administrativos.

Actualmente son fundamentales debido a que automatizan procesos y optimizan recursos.

\section{Desarrollo Web}

El desarrollo web consiste en crear aplicaciones accesibles desde navegadores mediante tecnologías frontend y backend.

Las aplicaciones web modernas permiten acceso multiplataforma y facilidad de mantenimiento.

\section{Bases de Datos}

Las bases de datos permiten almacenar información estructurada de manera organizada y segura.

En este proyecto se utiliza MySQL debido a su estabilidad y compatibilidad con PHP.

\section{Docker}

Docker es una plataforma basada en contenedores utilizada para crear entornos consistentes de desarrollo y producción.

\section{Git}

Git es un sistema de control de versiones distribuido utilizado para administrar cambios en proyectos de software.

% =========================================
% ALCANCE
% =========================================
\chapter{Alcance del Sistema}

\section{Funciones que realizará el sistema}

\begin{itemize}
    \item Registro de usuarios.
    
    \item Inicio de sesión.
    
    \item Administración de incidencias.
    
    \item Gestión de estados.
    
    \item Panel administrativo.
    
    \item Estadísticas.
\end{itemize}

\section{Funciones que no realizará}

\begin{itemize}
    \item Reparaciones automáticas.
    
    \item Gestión financiera.
    
    \item Integración con servicios externos.
\end{itemize}

% =========================================
% REQUERIMIENTOS
% =========================================
\chapter{Requerimientos del Sistema}

\section{Requerimientos Funcionales}

\begin{itemize}
    \item Registro de usuarios.
    
    \item Inicio de sesión.
    
    \item CRUD de incidencias.
    
    \item Administración de estados.
    
    \item Gestión de usuarios.
\end{itemize}

\section{Requerimientos No Funcionales}

\begin{itemize}
    \item Seguridad.
    
    \item Escalabilidad.
    
    \item Disponibilidad.
    
    \item Compatibilidad web.
\end{itemize}

% =========================================
% CASOS DE USO
% =========================================
\chapter{Casos de Uso}

\section{Registrar Incidencia}

\textbf{Actor principal:} Usuario.

\textbf{Descripción:}

El usuario inicia sesión y registra una incidencia proporcionando título, descripción y prioridad.

\textbf{Resultado esperado:}

La incidencia queda almacenada para seguimiento administrativo.

% =========================================
% HISTORIAS DE USUARIO
% =========================================
\chapter{Historias de Usuario}

\begin{itemize}
    \item Como alumno quiero registrar incidencias para reportar problemas institucionales.
    
    \item Como administrador quiero gestionar estados de incidencias.
    
    \item Como usuario quiero consultar mis reportes registrados.
\end{itemize}

% =========================================
% TECNOLOGÍAS
% =========================================
\chapter{Tecnologías Seleccionadas}

\begin{longtable}{|p{5cm}|p{7cm}|}
\hline
\textbf{Tecnología} & \textbf{Uso} \\
\hline
PHP 8 & Backend \\
\hline
MySQL & Base de datos \\
\hline
Docker & Contenedores \\
\hline
HTML5 & Frontend \\
\hline
CSS3 & Diseño visual \\
\hline
JavaScript & Interactividad \\
\hline
GitHub & Control de versiones \\
\hline
Chart.js & Estadísticas \\
\hline
\end{longtable}

% =========================================
% ARQUITECTURA
% =========================================
\chapter{Arquitectura del Sistema}

El sistema utiliza una arquitectura cliente-servidor.

\section{Cliente}

Representa la interfaz web accesible desde navegadores.

\section{Servidor}

Implementado mediante Apache y PHP.

\section{Base de Datos}

MySQL almacena toda la información relacionada con usuarios e incidencias.

% =========================================
% ESTRUCTURA DEL PROYECTO
% =========================================
\chapter{Estructura del Proyecto}

\begin{lstlisting}
project/
│
├── config/
├── database/
├── public/
├── src/
├── templates/
├── Dockerfile
└── docker-compose.yml
\end{lstlisting}

% Detalles añadidos: descripción por carpeta
\paragraph{Notas:}
\begin{itemize}
    \item \texttt{config/} contiene \texttt{database.php} con la conexión mysqli y el ajuste de charset a \texttt{utf8mb4}.
    \item \texttt{src/} contiene \texttt{auth.php} con helpers (find\_user\_by\_email, register\_user, ensure\_logged\_in, current\_user y fix\_encoding).
    \item \texttt{public/} es el DocumentRoot montado en el contenedor web (\texttt{/var/www/html/public}) y aloja las páginas públicas (login, register, dashboard, CRUD de incidencias).
    \item \texttt{database/} aloja el esquema (schema.sql), los seeds (seed.sql) y scripts utilitarios de conversión/limpieza.
\end{itemize}

\chapter{Detalles de implementación}
En esta sección se documentan piezas concretas del código que conviene conocer al ampliar o mantener la aplicación.

\section{Conexión a la base de datos}
El archivo principal de conexión es \texttt{config/database.php}. Contiene una conexión mysqli simple que apunta al host \texttt{mysql} (nombre del servicio Docker) y selecciona la base de datos \texttt{fesaragon}. Además fuerza el conjunto de caracteres UTF-8 extendido:
\begin{lstlisting}[language=PHP]
$conexion = mysqli_connect('mysql','root','root','fesaragon');
if(!$conexion) die('Error de conexion');
mysqli_set_charset($conexion, 'utf8mb4');
\end{lstlisting}

\section{Funciones de autenticación y utilidades}
Archivo: \texttt{src/auth.php}. Las funciones relevantes que se usan en las vistas son:
\begin{itemize}
    \item \texttt{find\_user\_by\_email(\$email)} — retorna el registro de \texttt{usuarios} por correo.
    \item \texttt{register\_user(\$nombre,\$correo,\$cuenta,\$password)} — registra un usuario usando \texttt{password\_hash} para guardar la contraseña.
    \item \texttt{ensure\_logged\_in()} — redirige al login si no hay sesión activa.
    \item \texttt{current\_user()} — devuelve los datos del usuario actual según \texttt{\$_SESSION['user\_id']}.
    \item \texttt{fix\_encoding(\$s)} — función práctica para corregir mojibake frecuente en algunos campos (reemplaza secuencias como \texttt{"á"} por \texttt{"á"}).
\end{itemize}

Ejemplo de uso en \texttt{login.php}:
\begin{lstlisting}[language=PHP]
$u = find_user_by_email(
        mysqli_real_escape_string($conexion, $correo)
);
if($u && password_verify($password, $u['password'])){
    // login OK: setear session y redirigir
}
\end{lstlisting}

\section{Manejo de encoding y mojibake}
La aplicación contiene utilitarios para intentar arreglar casos comunes de mojibake al mostrar datos; sin embargo, la estrategia correcta consiste en convertir la base de datos a \texttt{utf8mb4} y, si es necesario, ejecutar scripts de reparación sobre columnas afectadas (ver capítulo de migraciones).

\chapter{Uso de Contenedores y despliegue}
El proyecto se despliega con Docker Compose. El fichero \texttt{docker-compose.yml} define dos servicios principales: \texttt{web} y \texttt{mysql}.

\section{Resumen de servicios}
\begin{itemize}
    \item \texttt{web}: construye la imagen a partir del Dockerfile del repo, se monta la carpeta raíz del proyecto en \texttt{/var/www/html} dentro del contenedor y expone el puerto 80 del contenedor en el puerto 8080 del host. Nombre del contenedor: \texttt{fes\_web}.
    \item \texttt{mysql}: usa la imagen oficial \texttt{mysql:8}, crea una base de datos \texttt{fesaragon} y utiliza contraseña root \texttt{root} en el entorno de desarrollo. Nombre del contenedor: \texttt{fes\_mysql}.
\end{itemize}
\newpage
\section{Extracto relevante de \texttt{docker-compose.yml}}
\begin{lstlisting}
services:
    web:
        build: .
        container_name: fes_web
        ports:
            - "8080:80"
        volumes:
            - ./:/var/www/html
        depends_on:
            - mysql

    mysql:
        image: mysql:8
        container_name: fes_mysql
        environment:
            MYSQL_ROOT_PASSWORD: root
            MYSQL_DATABASE: fesaragon
        ports:
            - "3307:3306"
        volumes:
            - ./database/schema.sql:/docker-entrypoint-initdb.d/1-schema.sql
            - ./database/seed.sql:/docker-entrypoint-initdb.d/2-seed.sql
\end{lstlisting}

\section{Comandos recomendados (Windows PowerShell)}
\begin{verbatim}
# Construir y levantar servicios
docker-compose build --no-cache
docker-compose up -d

# Ver logs del web
docker-compose logs -f web

# Para detener y remover contenedores y volúmenes locales (destrucción de datos)
docker-compose down -v
\end{verbatim}

\chapter{Modelo de Base de Datos}
Se presenta aquí el DDL principal extraído de \texttt{database/schema.sql} con comentarios.

\section{Tabla Usuarios}

\begin{itemize}
    \item id
    
    \item nombre
    
    \item correo
    
    \item password
    
    \item rol
\end{itemize}

\section{Tabla Incidencias}

\begin{itemize}
    \item id
    
    \item titulo
    
    \item descripcion
    
    \item estado
    
    \item prioridad
\end{itemize}

\newpage
\section{Diagrama Entidad--Relación}
Para facilitar la comprensión de las relaciones entre tablas, a continuación se incluye un diagrama entidad-relación compacto que muestra entidades, claves primarias, claves foráneas y cardinalidades.


\begin{figure}[H]
\centering
\resizebox{\textwidth}{!}{
\begin{tikzpicture}[
    node distance=3cm,
    every node/.style={font=\small}
]

% =========================
% ENTIDADES
% =========================

\node[entity] (usuarios)
{
    \shortstack{
        \textbf{usuarios} \\
        \small id (PK) \\
        nombre \\
        correo \\
        cuenta \\
        password \\
        rol
    }
};

\node[entity, right=5cm of usuarios] (incidencias)
{
    \shortstack{
        \textbf{incidencias} \\
        \small id (PK) \\
        titulo \\
        descripcion \\
        tipo \\
        prioridad \\
        estado \\
        created\_at \\
        updated\_at
    }
};

\node[entity, above=3cm of incidencias] (edificios)
{
    \shortstack{
        \textbf{edificios} \\
        \small id (PK) \\
        nombre \\
        pisos
    }
};

\node[entity, right=5cm of edificios] (salones)
{
    \shortstack{
        \textbf{salones} \\
        \small id (PK) \\
        nombre \\
        edificio\_id (FK)
    }
};

\node[entity, below=3cm of salones] (zonas)
{
    \shortstack{
        \textbf{zonas} \\
        \small id (PK) \\
        nombre
    }
};

% =========================
% RELACIONES
% =========================

\draw[line]
(usuarios.east) -- node[above]{0..*}
(incidencias.west);

\draw[line]
(incidencias.north) -- node[left]{0..1}
(edificios.south);

\draw[line]
(edificios.east) -- node[above]{1..*}
(salones.west);

\draw[line]
(incidencias.east) -- node[right]{0..1}
(salones.south west);

\draw[line]
(incidencias.south east) -- node[right]{0..1}
(zonas.west);

\end{tikzpicture}
}

\caption{Diagrama entidad-relación del sistema}
\label{fig:er-diagrama}
\end{figure}





\noindent Documento actualizado con información extraída directamente del árbol del proyecto para que el reporte refleje la implementación real.

% =========================================
% CONTROL DE VERSIONES
% =========================================
\chapter{Estrategia de Control de Versiones}

GitHub fue utilizado para administrar el código fuente del proyecto.

El proyecto implementa:

\begin{itemize}
    \item Commits.
    
    \item Ramas.
    
    \item Pull Requests.
    
    \item Trabajo colaborativo.
\end{itemize}

% =========================================
% SEGURIDAD
% =========================================
\chapter{Seguridad del Sistema}

El sistema implementa medidas de seguridad importantes como:

\begin{itemize}
    \item Contraseñas cifradas.
    
    \item Validación de sesiones.
    
    \item Prepared Statements.
    
    \item Protección contra SQL Injection.
    
    \item Restricción por roles.
\end{itemize}

% =========================================
% PRUEBAS
% =========================================
\chapter{Estrategia de Pruebas}

\section{Pruebas Funcionales}

\begin{itemize}
    \item Registro de usuarios.
    
    \item Inicio de sesión.
    
    \item Registro de incidencias.
    
    \item Cambio de estados.
\end{itemize}

\section{Pruebas de Carga}

Se realizaron pruebas utilizando herramientas como k6 y ApacheBench.

% =========================================
% RIESGOS
% =========================================
\chapter{Riesgos del Proyecto}

\begin{itemize}
    \item Fallos del servidor.
    
    \item Vulnerabilidades de seguridad.
    
    \item Pérdida de información.
    
    \item Problemas de conectividad.
\end{itemize}

% =========================================
% CONCLUSIONES
% =========================================
\chapter{Conclusiones}

El sistema desarrollado representa una solución eficiente para optimizar la administración de incidencias dentro de la institución.

La plataforma permite centralizar información, mejorar la trazabilidad y reducir tiempos de atención.

Además, el uso de tecnologías modernas garantiza estabilidad, escalabilidad y facilidad de mantenimiento.

% =========================================
% REFERENCIAS
% =========================================
\chapter{Referencias}

\begin{itemize}

    \item PHP Documentation. Recuperado de: \\
    \url{https://www.php.net/docs.php}

    \item Docker Documentation. Recuperado de: \\
    \url{https://docs.docker.com/}

    \item MySQL Documentation. Recuperado de: \\
    \url{https://dev.mysql.com/doc/}

    \item Git Documentation. Recuperado de: \\
    \url{https://git-scm.com/doc}

    \item MDN Web Docs - HTML, CSS y JavaScript. Recuperado de: \\
    \url{https://developer.mozilla.org/}

    \item Chart.js Documentation. Recuperado de: \\
    \url{https://www.chartjs.org/docs/latest/}

    \item Apache HTTP Server Documentation. Recuperado de: \\
    \url{https://httpd.apache.org/docs/}

    \item GitHub Documentation. Recuperado de: \\
    \url{https://docs.github.com/}

    \item Repositorio oficial del proyecto. Recuperado de: \\
    \url{https://github.com/raynierjetsahe320/VinculacionEmpresarial}

\end{itemize}

\end{document}